% NeurIPS 2026 typography; the six proposal sections follow ECSE 552 requirements
% exactly, with paragraph subtitles mapping to items A1--A10 (marked in comments).
% Limit: about two pages, under 1,000 words, plus one page for references.
% Draft for team review (10 Oct 2026). Red (\tbd) = values the team must agree on;
% \review = text still to review (yellow in Word, plain in the PDF). Revise the prose freely; it is a starting point.
% Literature notes: docs/research/action-free-robot-learning/
\documentclass{article}
\PassOptionsToPackage{numbers,sort&compress}{natbib}
\usepackage[preprint]{neurips_2026}
% Course document, not a preprint: drop the "Preprint." notice on page 1.
\makeatletter\renewcommand{\@noticestring}{}\makeatother
\usepackage[utf8]{inputenc}
\usepackage[T1]{fontenc}
\usepackage{amsmath}
\usepackage{amssymb}
\usepackage{microtype}
\usepackage{xcolor}
\usepackage{booktabs}
\usepackage{array}
\usepackage[hidelinks]{hyperref}

% Tighter section and paragraph spacing than NeurIPS defaults, to fit two pages.
% Wrapped in \AtBeginDocument so pandoc (make docx) does not expand it.
\makeatletter
\AtBeginDocument{%
  \renewcommand{\section}{\@startsection{section}{1}{\z@}%
    {-1.4ex \@plus -0.3ex \@minus -0.2ex}{0.6ex \@plus 0.2ex}{\large\bf\raggedright}}%
  \renewcommand{\paragraph}{\@startsection{paragraph}{4}{\z@}%
    {0.6ex \@plus 0.3ex \@minus 0.2ex}{-1em}{\normalsize\bf}}%
  \setlength{\parskip}{3pt}}
\makeatother

\newcommand{\todo}[1]{\textcolor{red}{[TO WRITE: #1]}}
\newcommand{\decide}[1]{\textcolor{red}{[TEAM: #1]}}
% Values the team still has to agree on (red) and member names to fill (light red).
\newcommand{\tbd}[1]{\textcolor{red}{#1}}
\definecolor{teamlight}{HTML}{E06666}
\newcommand{\member}[1]{\textcolor{teamlight}{#1}}
% Text the team still has to review: unchanged in the PDF, highlighted yellow in
% Word. Wrap one paragraph at a time (\textcolor cannot span paragraphs).
\definecolor{review}{HTML}{000000}
\newcommand{\review}[1]{\textcolor{review}{#1}}
% Starting ideas, rendered in red so they are never submitted by accident.
\newenvironment{ideas}
  {\color{red}\small\begin{itemize}\setlength{\itemsep}{0pt}}
  {\end{itemize}}

\title{Learning to Manipulate a New Object\\from Action-Free Videos}
\author{
  Group 9 \\
  McGill University \\
  ECSE 552 --- Project proposal
}

\begin{document}
\maketitle

\section*{Background}
% A1
\paragraph{Application.}
We teach a robot arm that manipulates one object to handle new ones from videos
alone: warehouse and home robots meet new items constantly, and recording paired
images and motor commands for each requires operating the robot. We study this in
OCBench simulation (UR5e arm, Robotiq gripper) \citep{ocbench2026}.

% A2
\paragraph{Problem and need.}
Visuomotor policies imitate images paired with motor commands, and such pairs are
costly: 76k teleoperated demonstrations took 50 people a year
\citep{khazatsky2024droid}. Videos without commands are cheap (e.g., generated by
world models or video generation models \citep{jang2025dreamgen}). An inverse
dynamics model (IDM), trained on few labeled demonstrations
\citep{morin2026idmssil}, infers the command between consecutive frames and
labels such videos for imitation \citep{torabi2018bco,baker2022vpt}. Here the target videos show new objects, unlike the IDM's training data. Formally, given labeled cube transitions
$\mathcal{D}_A=\{(o_t,o_{t+1},a_t)\}$ and unlabeled target videos
$\mathcal{D}_B=\{o_t\}$, we seek an image-only policy $\pi(a \mid o)$ that
succeeds on the new objects under R1--R6.

% A3
\paragraph{Constraints.}
\tbd{Red values are still to be agreed.}

\begin{center}
\small
\begin{tabular}{@{}l>{\raggedright\arraybackslash}p{2.0cm}>{\raggedright\arraybackslash}p{7.6cm}>{\raggedright\arraybackslash}p{2.3cm}@{}}
\toprule
ID & Requirement & Definition & Target \\
\midrule
R1 & Supervision & Target data are RGB frames only & Zero target labels \\
R2 & Label budget & Labeled cube demonstrations, identical for every method & \tbd{300} demos \\
R3 & Inputs & No simulator state in any learned model & Images only \\
R4 & Target & Success gain over the cube-only policy on each target, as a fraction of the oracle's gain & $\geq$ \tbd{50\%} of gap \\
R5 & Retention & Drop in cube success after adding target data & $\leq$ \tbd{10} points \\
R6 & Real time & Policy inference per call, within one 50\,Hz control step (IDM offline) & $<$ \tbd{20}\,ms \\
\bottomrule
\end{tabular}
\end{center}

\section*{Aims/Goals}
% A3
\paragraph{Goals.}
Meet R4 and R5 on each new object with an image-only policy trained on cube
demonstrations and action-free target videos; R1--R3 hold by design and R6 is
measured.

% A4
\paragraph{Scope.}
We use OCBench's visual single-block task; the source is its default
$6 \times 6 \times 6$\,cm red cube. \review{The targets separate the factors: (A) a
$6 \times 6 \times 6$\,cm blue cube (appearance only), (B) a
$9 \times 6 \times 6$\,cm red block (geometry, same grasp width), (C) a
$9 \times 7 \times 6$\,cm red block (geometry and grasp width), and, optionally,
(D) a $9 \times 7 \times 6$\,cm blue block (all combined).} Expert grasps are
adapted to each; robot, cameras, lighting, and task stay fixed. Human videos,
real hardware, and large pretrained models are out of scope.

% A5
\paragraph{Proposed solution.}
We adapt the IDM to each target before labeling its videos. The IDM sees frame differences, which emphasize the unchanged arm, and
domain-adversarial training \citep{ganin2016dann} makes its features
indistinguishable between cube and target frames while still predicting cube
actions. One policy per target then trains on cube demonstrations plus pseudo-labeled
target videos. \review{Adaptation should help most on A, whose
actions match the cube's.}

\section*{Methodology}
% A5, A6
\paragraph{Design and implementation.}
\review{Observations are $224{\times}224$ images from the front, side, and wrist
cameras; actions are seven normalized joint and gripper increments at 50\,Hz. The
IDM's convolutional encoder, shared across cameras, embeds $o_t$, $o_{t+1}$, and
their difference, and an MLP regresses $a_t$ with mean squared error. A
discriminator classifies these features as cube or target (binary
cross-entropy); a gradient reversal layer returns its gradient to the encoder
with flipped sign, scaled by a weight $\lambda$ that ramps up as in DANN;
$\lambda=0$ gives the unadapted IDM. The policy shares the encoder, takes $o_t$,
and is trained with flow matching to predict the next \tbd{25} actions, executed
open loop as in OCBench \citep{ocbench2026}. All methods share encoder, data, and
budget (PyTorch).}

% A7
\paragraph{Validation, testing, and evaluation.}
\review{Splits are made by episode. Target actions never select models: the
$\lambda$ schedule is fixed in advance, IDM checkpoints follow source validation
error, and policies train for a fixed number of steps. We report target success,
cube retention, per-component IDM error, and latency over three seeds and 100
episodes per method and object, with identical initial states.}

% A8
\paragraph{Alternatives.}
\review{We compare (a) a cube-only behavioral cloning policy, (b) an unadapted
IDM, (c) our adversarial IDM, and (d) an oracle trained with true target actions;
optionally, (e) applies the alignment to the policy encoder instead. (b) versus
(c) isolates adaptation; (a) and (d) anchor R4.}

\section*{Feasibility and resources}
% A9
\paragraph{Dataset, resources, and plan.}
We generate all data in OCBench with MuJoCo: \tbd{300} labeled cube demonstrations
and \tbd{200} action-free videos per target, keeping true target actions only for
evaluation and the oracle; no licensing or ethics constraints apply. Training
uses one NVIDIA A10G GPU; the pilot estimates compute. \review{A pilot by
\tbd{October 21} checks expert success and, on the cube and each target, the
cube-only policy's success and the IDM's error.} Core methods follow by
\tbd{November 1}; evaluation by \tbd{November 18}.

\section*{Risk analysis}
% A10
\paragraph{Pitfalls and mitigation.}
 \review{Adversarial training aligns feature distributions, not the
action each feature maps to \citep{zhao2019invariant}; we expect actions to stay
readable from arm motion, which per-component errors will test.} If the expert
fails on a target, we adjust the object or grasp\review{; if the cube-only policy
fails on the cube, we use OCBench's shorter lite variant}. If alignment hurts, we report negative transfer; optional experiments (D, e) go
first if time runs short.

\section*{Role of group members}
All members collaborate and review the report.

\begin{center}
\small
\begin{tabular}{@{}>{\raggedright\arraybackslash}p{1.9cm}>{\raggedright\arraybackslash}p{3.4cm}>{\raggedright\arraybackslash}p{7.8cm}@{}}
\toprule
Member & Role & Tasks \\
\midrule
\member{MEMBER 1} & Data and evaluation & Simulation, expert demos, datasets, splits, rollouts, metrics \\
\member{MEMBER 2} & Adversarial adaptation & Discriminator, adversarial training, tuning, experiments \\
\member{MEMBER 3} & Inverse dynamics model & IDM architecture and training, pseudo-labeling, error analysis \\
\member{MEMBER 4} & Baselines and alternatives & Policy training, baselines, alternative methods, comparisons \\
\bottomrule
\end{tabular}
\end{center}

% Only works cited with \citep{} appear in the references.
\clearpage
{\footnotesize\textbf{Generative-AI disclosure.} Codex (OpenAI) and Claude Code
(Anthropic) assisted with literature search, bibliography formatting, document
structure, and prose. We used them to survey recent work and to format the
document; the team revised all text and verified every claim and citation against
the original sources.}

{\small
\bibliographystyle{unsrtnat}
\bibliography{../references}}
\end{document}
